\begin{lemma}
\label{lemma-flat-factors-free}
Let $M$ be an $R$-module.  The following are equivalent:
\begin{enumerate}
\item $M$ is flat.
\item If $f: R^n \to M$ is a module map and $x \in \Ker(f)$, then there
are module maps $h: R^n \to R^m$ and $g: R^m \to M$ such that
$f = g \circ h$ and $x \in \Ker(h)$.
\item Suppose $f: R^n \to M$ is a module map, $N \subset \Ker(f)$ any
submodule, and $h: R^n \to R^{m}$ a map such that $N \subset \Ker(h)$
and $f$ factors through $h$.  Then given any $x \in \Ker(f)$ we can find a map
$h': R^n \to R^{m'}$ such that $N + Rx \subset \Ker(h')$ and $f$
factors through $h'$.
\item If $f: R^n \to M$ is a module map and $N \subset \Ker(f)$ is a
finitely generated submodule, then there are module maps $h: R^n \to
R^m$ and $g: R^m \to M$ such that $f = g \circ h$ and $N \subset
\Ker(h)$.
\end{enumerate}
\end{lemma}

\begin{proof}
That (1) is equivalent to (2) is just a reformulation of the equational
criterion for flatness\footnote{In fact, a module map $f : R^n \to M$
corresponds to a choice of elements $x_1, x_2, \ldots, x_n$ of $M$
(namely, the images of the standard basis elements $e_1, e_2, \ldots,
e_n$); furthermore, an element $x \in \Ker(f)$ corresponds to a
relation between these $x_1, x_2, \ldots, x_n$ (namely, the relation
$\sum_i f_i x_i = 0$, where the $f_i$ are the coordinates of $x$).
The module map $h$ (represented as an $m \times n$-matrix)
corresponds to the matrix $(a_{ij})$ from Lemma \ref{lemma-flat-eq},
and the $y_j$ of Lemma \ref{lemma-flat-eq} are the images of the
standard basis vectors of $R^m$ under $g$.}.   To show
(2) implies (3), let $g: R^m \to M$
be the map such that $f$ factors as $f = g \circ h$.  By (2) find $h'': R^m
\to R^{m'}$ such that $h''$ kills $h(x)$ and $g: R^m \to M$
factors through $h''$.  Then taking $h' = h'' \circ h$ works.  (3) implies (4)
by induction on the number of generators of $N \subset \Ker(f)$ in (4).
Clearly (4) implies (2).
\end{proof}
